%%%%%     PACKS     %%%%%
\documentclass[12pt]{article}
\usepackage[margin=1in,headsep=.60in]{geometry}
\usepackage[utf8]{inputenc}
\usepackage[table, dvipsnames]{xcolor}
\usepackage{array}
\usepackage{amsmath}
\usepackage{booktabs}
\usepackage{mdframed} %For box around text
\usepackage{hyperref} %For hyperlinks in the PDF
\hypersetup{
    colorlinks=true,
    linkcolor=blue,
    filecolor=magenta,
    urlcolor=cyan,
}
\urlstyle{same}
\usepackage{amsthm}
\usepackage{amssymb}
\usepackage{amsfonts}
\usepackage{siunitx}
\usepackage{graphicx}
\usepackage{caption}
\usepackage{pgfplots}
\graphicspath{{Images/}}
\usepackage[colorinlistoftodos]{todonotes}
\usepackage{cleveref}
\usepackage[labelformat=simple]{subcaption}
\usepackage{grffile}
\usepackage{gensymb}
\usepackage{float}
\usepackage[shortlabels]{enumitem}
\usepackage{enumitem}
\setlistdepth{9}
\usepackage{biblatex} %Imports biblatex package
\addbibresource{references.bib}
\usepackage{outlines}
\usepackage{minted}
\usepackage{pdflscape}
\usepackage{everypage}
\usepackage{multirow}
\usepackage{multicol}
\usepackage{afterpage}
\usepackage[labelfont=bf, font=small]{caption}
\usepackage{upgreek}
\usepackage{tikz}
\usepackage{xcolor}
\usepackage{listings}

% style sheet for code
\definecolor{codegreen}{rgb}{0,0.6,0}
\definecolor{codegray}{rgb}{0.5,0.5,0.5}
\definecolor{codepurple}{rgb}{0.58,0,0.82}
\definecolor{backcolour}{rgb}{0.95,0.95,0.92}

\lstdefinestyle{mystyle}{
    backgroundcolor=\color{backcolour},   
    commentstyle=\color{codegreen},
    keywordstyle=\color{magenta},
    numberstyle=\tiny\color{codegray},
    stringstyle=\color{codepurple},
    basicstyle=\ttfamily\footnotesize,
    breakatwhitespace=false,         
    breaklines=true,                 
    captionpos=b,                    
    keepspaces=true,                 
    numbers=left,                    
    numbersep=5pt,                  
    showspaces=false,                
    showstringspaces=false,
    showtabs=false,                  
    tabsize=2
}

\lstset{style=mystyle}

\usepackage{lastpage}
%%%%%     COMMANDS     %%%%%
%%%% Blue box for subsection text (not figures) %%%%
\newenvironment{bluebox}
  {\begin{mdframed}[backgroundcolor=blue!5,linecolor=blue!40,roundcorner=8pt]}
  {\end{mdframed}}

%%%% Simple neutral box for figures %%%%
\newenvironment{figbox}
  {\begin{mdframed}[roundcorner=8pt,shadow=true,shadowsize=4pt,shadowcolor=black!40]}
  {\end{mdframed}}

\begin{document}

\begin{titlepage}
\newcommand{\HRule}{\rule{\linewidth}{0.5mm}}

\center

\textsc{\LARGE Aarhus University}\\[1.5cm]
\textsc{\Large Department Of Mathematics}\\[0.5cm]
\textsc{\large Introduction to Algebra 1}\\[0.5cm]
    

\HRule\\[0.4cm]
    \center 
    {\huge\bfseries  Handin 4}\\[0.4cm] % Title of your document
\HRule\\[1.5cm]

\begin{minipage}{0.4\textwidth}
        \begin{flushleft}
            \large
            \textit{Author}\\   
            Søren M. \textsc{Damsgaard}\\
                 % Your name
        \end{flushleft}
    \end{minipage}
~
    \begin{minipage}{0.4\textwidth}
        \begin{flushright}
            \large
            \textit{Student number}\\
            \textbf{202309814}\\
                 % Studienummer\
            
        \end{flushright}
    \end{minipage}

\vfill\vfill\vfill % Position the date 3/4 down the remaining page
    
    {\large\today}

\vfill\vfill
    \includegraphics[width=0.2\textwidth]{Aarhus_University_seal.png}\\[1cm] % Include a department/university logo - this will require the graphicx package
    

\vfill
\end{titlepage}
%%%%%      CHAPTERS      %%%%%
\hspace{0.02cm}
\section*{Non-Goofy Math IV}
\#ConsistencyIsForFleshBags

%Note that a prime factor is a factor of a number that is a prime,
\section{Show that a natural number congruent to 3 modulo 4 has a prime factor congruent to 3 modulo 4.}
Let $n \in \mathbb{N}$ and $n \equiv 3 \mod 4$ then per definition 9.1~\cite{Algebra} then it holds that
\[
4 | n-3
\]
and per the defintion of divisibility, then it holds for some $q \in \mathbb{W}$ that
\[
n-3 = 4 \cdot q \rightarrow n = 4q + 3
\]
and because of $+3$ it can't satisfy the formula for being even $n = 2x$ for $x \in \mathbb{N}$ so it holds that $n$ must be odd
\newline
and from defintion 4.2~\cite{Algebra} it holds that $n$ has a prime factorization, and that they must be odd. 
\newline
\newline
For modulo 4 there are four congruent classes $[0]$, $[1]$, $[2]$ and $[3]$, and we need to show that odd numbers are either in congruent class $[1]$ or $[3]$ and even numbers are in $[0]$ or $[4]$.\\
To do that we define a number $k$
\begin{align*}
    &\text{if } k \equiv 0 \mod 4, \text{ then }k = 4x + 0 = 2(2x)\\
    &\text{if } k \equiv 2 \mod 4, \text{ then }k = 4x + 2 = 2(2x+1)
\end{align*}
In both cases, $k$ satisfy the evening ($k$ is even) and in the other two cases for $[1]$ and $[3]$ then $k$ must be odd.\\
This means any odd prime factor must be either $k \equiv 1$ or $k \equiv 4$\\
\newline
Now assume, for proof by contradiction, that every odd prime is congruent 1 modulo 4:\\
and then we look at lemma 9.2~\cite{Algebra} where it holds for any natural number (also primes), if they are congruent to the same modulo, then their product is alos congruent to the same modulo.\\
So if every odd prime is congruent 1 modulo 4, then their product is also congruent 1 modulo 4, but this contradicts $n$ which is a product of odd primes and is congruent 3 modulo 4, this contradictions means the original assumption
of every odd prime being congruenr 1 modulo 4 is false, so there must be atleast one prime factor of $n$ that is congruent 3 modulo 4.


\section{Show that there are infinitely many prime numbers that are congruent to 3 modulo 4}
To show this, we will use Euclid's proof 5.1~\cite{Algebra} and the result from the previous section, via a proof by contradiction.\\
\newline
Assume $p_1, \hdots , p_N$ to be a discrete list of all prime numbers that are congruent 3 modulo 4, then construct a new number $m$ such
\[
m = 4\cdot p_1 \cdot p_2 \cdot \hdots \cdot p_N
\]
that it is congruent 0 modulo 4, per the definition of division $4|m$ then $m = 4 \cdot q$ where $q = p_1 \cdot p_2 \cdot \hdots \cdot p_N$.\\
and if one subtracts 1 from $m$ then we get
\[
m = 4p_1p_2 \hdots p_N - 1 \equiv -1 \equiv 3 \mod 4
\]
and according to theorem 4.3~\cite{Algebra} since $m$ is a natural number bigger than 1, then it holds that it is divisible by one of it's prime factors, 
but $m$ isn't divisible by any of the primes $p_1, \hdots, p_N$, since that would mean any of them could divide into 1, which they can't, this contradiction means the original assumption must be wrong and the opposite is true and there is a prime factor $p_{new}$
 of $m$ that isn't in the original list, further, according to the previous result from section 1, $m$ has atleast one prime factor that is congruent 3 modulo 4 and we let this be  $p_{new} \equiv 3 \mod 4$, so there must be an infinite amount of primes that are congruent 3 modulo 4 
%Bro, there are infinite primes, so there is infinite odd primes, so there must be infinite primes congruent 3 modulo 4, checkmate
%Broooo....you're so right
\printbibliography
\end{document}